\section{Yêu cầu phi chức năng}
\label{sec:yeucau_phichucnang}

Các yêu cầu phi chức năng xác định những mục tiêu chất lượng cần xem xét trong quá trình thiết kế và hiện thực MyHCMUT Mobile. Mức độ đáp ứng được đánh giá thông qua kết quả kiểm thử và các giới hạn trình bày tại Chương~6.

\begin{table}[H]
\centering\small
\caption{Yêu cầu phi chức năng của MyHCMUT Mobile}
\label{tab:nfr}
\begin{tabularx}{\textwidth}{|L{1.8cm}|L{3.4cm}|Y|}
\hline
\textbf{Mã} & \textbf{Yêu cầu} & \textbf{Nội dung đặc tả} \\
\hline
NFR-01 & Tính ổn định & Cần xử lý mất kết nối, quá thời gian chờ và lỗi hệ thống nguồn. Cần phân biệt dữ liệu rỗng với dữ liệu chưa tải được; chỉ báo thao tác thành công sau xác nhận của backend. \\
\hline
NFR-02 & Khả năng sử dụng & Giao diện cần nhất quán, phù hợp thao tác di động; biểu mẫu nhiều thông tin được chia bước. Cần thể hiện trạng thái tải, lỗi, kết quả và điều kiện tiếp tục thao tác. \\
\hline
NFR-03 & Bảo mật và phân quyền & Cần xác thực trước khi truy cập nghiệp vụ. UI phản ánh quyền tài khoản; backend cần kiểm tra quyền theo nghiệp vụ, đối tượng và bước xử lý, kể cả yêu cầu gọi trực tiếp. \\
\hline
NFR-04 & Tính nhất quán dữ liệu & Dữ liệu chính thức do hệ thống nguồn quản lý; dữ liệu cục bộ chỉ hỗ trợ sử dụng. Sau thao tác ghi, cần phản ánh kết quả được backend ghi nhận; kết quả chưa xác định khi mất phản hồi không được coi là lưu thất bại chắc chắn. \\
\hline
NFR-05 & Khả năng bảo trì và mở rộng & Mục tiêu thiết kế là tổ chức mã theo nhóm nghiệp vụ và tách thành phần dùng chung, hạn chế phụ thuộc không cần thiết để hỗ trợ thay đổi và kiểm thử. \\
\hline
\end{tabularx}
\end{table}

Các yêu cầu trên được dùng làm tiêu chí đối chiếu khi đánh giá hệ thống. Kết quả đánh giá và các giới hạn tương ứng được trình bày tại Chương~6.

\section{Kết luận chương}
\label{sec:ketluan_chuong4}

Chương này đã xác định các nhóm người dùng, yêu cầu chức năng và yêu cầu phi chức năng của MyHCMUT Mobile. Các yêu cầu chức năng được phân tích theo những nhóm nghiệp vụ chính và được làm rõ thông qua các ca sử dụng, biểu đồ hoạt động và biểu đồ tuần tự đối với các quy trình quan trọng. Bên cạnh đó, các yêu cầu phi chức năng xác định những yêu cầu về tính ổn định, khả năng sử dụng, bảo mật và phân quyền, tính nhất quán dữ liệu, cũng như khả năng bảo trì và mở rộng của ứng dụng.

Các yêu cầu được xác định trong chương này là cơ sở cho việc thiết kế kiến trúc, tổ chức các thành phần ứng dụng và xây dựng các cơ chế tích hợp với hệ thống nghiệp vụ hiện hữu được trình bày trong Chương 5.
